\documentclass[12pt,a4paper]{article}
\usepackage[T1]{fontenc}
\usepackage{amsmath,amssymb,mathrsfs,tikz,times,pifont}
\usepackage{enumitem}
\usepackage{multicol}
\usepackage{lmodern}
\usetikzlibrary{trees}
\newcommand\circitem[1]{%
\tikz[baseline=(char.base)]{
\node[circle,draw=gray, fill=red!55,
minimum size=1.2em,inner sep=0] (char) {#1};}}
\newcommand\boxitem[1]{%
\tikz[baseline=(char.base)]{
\node[fill=cyan,
minimum size=1.2em,inner sep=0] (char) {#1};}}
\setlist[enumerate,1]{label=\protect\circitem{\arabic*}}
\setlist[enumerate,2]{label=\protect\boxitem{\alph*}}
%%%::::::by chnini ameur :::::::%%%
\everymath{\displaystyle}
\usepackage[left=1cm,right=1cm,top=1cm,bottom=1.7cm]{geometry}
\usepackage[colorlinks=true, linkcolor=blue, urlcolor=blue, citecolor=blue]{hyperref}
\usepackage{array,multirow}
\usepackage[most]{tcolorbox}
\usepackage{varwidth}
\usepackage{float} %pour utiliser l'option [H] qui force l'image à apparaître exactement à l'endroit où elle est placée dans le code.
\tcbuselibrary{skins,hooks}
\usetikzlibrary{patterns}
%%%::::::by chnini ameur :::::::%%%
\newtcolorbox{exa}[2][]{enhanced,breakable,before skip=2mm,after skip=5mm,
colback=yellow!20!white,colframe=black!20!blue,boxrule=0.5mm,
attach boxed title to top left ={xshift=0.6cm,yshift*=1mm-\tcboxedtitleheight},
fonttitle=\bfseries,
title={#2},#1,
% varwidth boxed title*=-3cm,
boxed title style={frame code={
\path[fill=tcbcolback!30!black]
([yshift=-1mm,xshift=-1mm]frame.north west)
arc[start angle=0,end angle=180,radius=1mm]
([yshift=-1mm,xshift=1mm]frame.north east)
arc[start angle=180,end angle=0,radius=1mm];
\path[left color=tcbcolback!60!black,right color = tcbcolback!60!black,
middle color = tcbcolback!80!black]
([xshift=-2mm]frame.north west) -- ([xshift=2mm]frame.north east)
[rounded corners=1mm]-- ([xshift=1mm,yshift=-1mm]frame.north east)
-- (frame.south east) -- (frame.south west)
-- ([xshift=-1mm,yshift=-1mm]frame.north west)
[sharp corners]-- cycle;
},interior engine=empty,
},interior style={top color=yellow!5}}
%%%%%%%%%%%%%%%%%%%%%%%
\usepackage{fancyhdr}
\usepackage{eso-pic}         % Pour ajouter des éléments en arrière-plan
% Commande pour ajouter du texte en arrière-plan
\usepackage{tkz-tab}
\usepackage{tikz}
\usepackage[siunitx, european, RPvoltages]{circuitikz}
\AddToShipoutPicture{
    \AtTextCenter{%
        \makebox[0pt]{\rotatebox{80}{\textcolor[gray]{0.7}{\fontsize{5cm}{5cm}\selectfont PGB}}}
    }
}
\usepackage{lastpage}
\fancyhf{}
\pagestyle{fancy}
\renewcommand{\footrulewidth}{1pt}
\renewcommand{\headrulewidth}{0pt}
\renewcommand{\footruleskip}{10pt}
\fancyfoot[R]{
\color{blue}\ding{45}\ \textbf{2026}
}
\fancyfoot[L]{
\color{blue}\ding{45}\ \textbf{Prof:M. BA}
}
\cfoot{\bf
\thepage /
\pageref{LastPage}}
% Création du compteur pour les exercices
\newcounter{question}
\renewcommand{\thequestion}{\arabic{question}}  % Définit l'affichage du compteur en chiffres arabes

% Définir la commande \exo
\newcommand{\exo}{\refstepcounter{question}\textbf{Question \thequestion} }
\begin{document}
\renewcommand{\arraystretch}{1.5}
\renewcommand{\arrayrulewidth}{1.2pt}
\hrule
\begin{center}
    \begin{tabular}{@{} p{5cm} p{5cm} p{5cm} @{}}
        Lycée Dindéfélo & \quad\quad\quad Test 04 & 07 Novembre 2026 \\
    \end{tabular}
    \\[-0.01cm]
    \hrule
\end{center}
\begin{center}
    \textbf{\Large Sommes Algébriques et Puissances de Radicaux} \\[0.2cm]
    \textbf{\large Professeur : M. BA} \\[0.2cm]
    \textbf{Classe : 3$^{\textbf{ème}}$M1} \\[0.2cm]
    \textbf{\small Durée : 10 minutes} \\[0.2cm]
    \textbf{\small Note :\quad\quad /5}
\end{center}

\textbf{\small Nom de l'élève :} \underline{\hspace{8cm}} \\[0.5cm]

\fbox{\textbf{\exo}} (1,5 point) Écrire les radicaux suivants sous la forme $a\sqrt{b}$ avec $a \in \mathbb{N}$ et $b \in \mathbb{N}$ le plus petit possible :
\[
\begin{aligned}
\sqrt{27}&=\sqrt{\dots \times 3}\\
&=\\
&=\\
&\\
\end{aligned}
\quad\quad\quad\quad\quad
\begin{aligned}
\sqrt{72}&=\sqrt{\dots\times\dots\times 2}\\
&=\\
&=\\
&=\\
\end{aligned}
\quad\quad\quad\quad\quad
\begin{aligned}
\sqrt{75}&=\sqrt{\dots\times 3}\\
&=\\
&=\\
&\\
\end{aligned}
\]

\vspace{0.8cm}

\fbox{\textbf{\exo}} (2 points) Simplifier les sommes algébriques suivantes :
\[
\begin{aligned}
A&= 2\sqrt{27} - \sqrt{75}\\
&=\\
&=\\
&=\\
&=\\
\end{aligned}
\quad\quad\quad\quad\quad\quad\quad\quad
\begin{aligned}
B&= \sqrt{72} + 3\sqrt{18} - \sqrt{50}\\
&=\\
&=\\
&=\\
&=\\
\end{aligned}
\]

\vspace{0.8cm}

\fbox{\textbf{\exo}} (1,5 point) Simplifier les expressions suivantes comportant des puissances sous le radical :
\[
\begin{aligned}
C&=\sqrt{5^6}\\
&=\\
&=\\
&=\\
\end{aligned}
\quad\quad\quad\quad\quad\quad\quad\quad
\begin{aligned}
D&=\sqrt{2^7}\\
&=\\
&=\\
&=\\
\end{aligned}
\]

\end{document}